\section{Introduction}

Intermediate representations (IRs) are at the heart of every modern compiler. These data structures
represent programs throughout compilation, connect individual compiler stages, and provide
abstractions to facilitate the implementation of analyses, optimizations, and program
transformations. A suitable IR highlights and exposes program properties that are important to the
transformations in a specific compiler stage. This reduces the complexity of optimizations and
simplifies their implementation.


Modern embedded systems have become increasingly parallel %
as system designers strive to improve their computational power and energy
efficiency. %
Increasing the number of cores in a system enables each core to be operated at a
lower clock frequency and supply voltage, improving
overall energy efficiency while providing sufficient system performance. %
Multi-core systems also reduce the total system cost by enabling the consolidation
of multiple functionalities onto a single chip.
In order to take full advantage of these systems, optimizing compilers need to
expose a program's available parallelism. %
This has led to an interest in developing more efficient program
representations~\cite{Cordes10,Lattner20}
and methodologies and frameworks~\cite{Castrillon13} for exposing the necessary information.

Data flow centric IRs, such as the Value (State) Dependence Graph
(V(S)DG)~\cite{Weise94, Johnson03, Johnson04, Lawrence07, Stanier211, Stanier12}, show promises as a
new class of IRs for optimizing compilers. %
These IRs are based on the observation that many optimizations require data flow
rather than control flow information, and shift the focus to explicitly expose
data instead of control flow. %
They represent programs in demand-dependence form, encode structured control
flow, and explicitly model data flow between operations. %
This raises the IR's abstraction level, permits simple and powerful
implementations of data flow optimizations, and helps to expose the inherent
parallelism in programs~\cite{Lawrence07, Johnson04, Stanier211}. %
However, the shift in focus from explicit control flow to only structured and
implicit control flow requires more sophisticated construction and destruction
methods~\cite{Weise94, Lawrence07, Stanier12}. %
In this context, Bahmann~\etal~\cite{Bahmann15} presents the \emph{Regionalized Value State
Dependence Graph} (RVSDG) and conclusively addresses the problem of intra-procedural control
flow recovery for demand-dependence graphs. %
They show that the RVSDG's restricted control flow constructs do not limit the complexity of
the recoverable control flow. %

In this work, we are concerned with the aspects of unified program
representation in the RVSDG. We present the required RVSDG constructs, consider
construction and destruction at the program level, and show feasibility and
practicality of this IR for optimizations by providing a practical compiler implementation.
Specifically, we make the following contributions: \emph{i)} A complete RVSDG specification,
including intra- and inter-procedural constructs. \emph{ii)} A complete description of RVSDG
construction and destruction, augmenting the previously proposed algorithms with the construction
and destruction of inter-procedural constructs, as well as the handling of intra-procedural
dependencies during construction. \emph{iii)} A presentation of Dead Node Elimination (DNE) and
Common Node Elimination (CNE) optimizations to demonstrate the RVSDG's utility. DNE combines dead
and unreachable code elimination, as well as dead function removal. CNE permits the removal of
redundant computations by detecting congruent operations. \emph{iv)} A publicly
available~\cite{jlm} prototype compiler that implements the discussed concepts.  It consumes and
produces LLVM IR, and is to our knowledge the first optimizing compiler that uses a demand
dependence graph as IR. \emph{v)} An evaluation of the RVSDG in terms of performance and size of
the produced code, as well as compile time and representational overhead.








Our results show that the RVSDG can serve as the IR in a compiler's optimization stage, producing
competitive code even with a conservative modeling of programs using a single memory state.
Even though this leaves significant parallelization potential
unused, it already yields satisfactory results compared to control-flow
based approaches.
This work paves the way for %
further exploration of the RVSDG's properties and their effect on optimizations
and analyses, as well as its usability in code generation for dataflow and
parallel architectures.





